\subsection{正函数的叠积分}

在本节其余部分，如无相反的明确说明，我们将用 X 表示一个局部紧空间，用 \(\Lambda: t \mapsto \lambda_{t}\) 表示一个 \(\mu\) -适宜映射（T 到 \(\mathcal{M}_{+}(X)\) 中的），并用 \(\nu\) 表示 \(\Lambda\) 的积分。

命题 3. --- 设 \(f\) 是定义在 X 上的一个数值函数 \(\geqslant 0\) 。a) 下列不等式成立：

\[
\int^ {*} f (x) d \nu (x) \geqslant \int^ {\bullet} d \mu (t) \int^ {*} f (x) d \lambda_ {t} (x) \geqslant \int^ {\bullet} d \mu (t) \int^ {\bullet} f (x) d \lambda_ {t} (x). \tag {6}
\]

\begin{enumerate}
\def\labelenumi{\alph{enumi})}
\setcounter{enumi}{1}
\tightlist
\item
  若 \(\Lambda\) 是淡连续的，则
\end{enumerate}

\[
\int^ {*} f (x) d \nu (x) \geqslant \int^ {*} d \mu (t) \int^ {*} f (x) d \lambda_ {t} (x).\tag{7}
\]

\begin{enumerate}
\def\labelenumi{\alph{enumi})}
\setcounter{enumi}{2}
\tightlist
\item
  若 \(\lambda_t^\bullet (1) <   + \infty\) 局部 \(\mu\) -几乎处处成立，则
\end{enumerate}

\[
\int^ {\bullet} f (x) d \nu (x) \geqslant \int^ {\bullet} d \mu (t) \int^ {*} f (x) d \lambda_ {t} (x) = \int^ {\bullet} d \mu (t) \int^ {\bullet} f (x) d \lambda_ {t} (x). \tag {8}
\]

设 \(g\) 是 \(X\) 上的一个下半连续函数，满足 \(f \leqslant g\) 。对每个 \(t \in T\) ，

\[
\int^ {*} f (x) d \lambda_ {t} (x) \leqslant \int^ {*} g (x) d \lambda_ {t} (x),
\]

因此，由 (4) 及 §1 的命题 4，

\[
\int^ {\bullet} d \mu (t) \int^ {*} f (x) d \lambda_ {t} (x) \leqslant \int^ {\bullet} d \mu (t) \int^ {*} g (x) d \lambda_ {t} (x) = \int^ {*} g (x) d \nu (x).
\]

于是，不等式 (6) 中的第一个由 \(\int^{*}f(x)d\nu(x)\) 的定义（Ch. IV, §1, No.~3, 定义 3）得出，第二个则由它立即得出。若 \(\Lambda\) 是淡连续的，则不等式 (7) 用 (5) 代替 (4)，以类似的方式得到证明。

现在转向 (8) 的证明。映射 \(t \mapsto \lambda_t^*(1)\) 是可测的，且局部 \(\mu\) -几乎处处有限。因此，集 \(\mathfrak{K}\) （由 \(T\) 中使 \(t \mapsto \lambda_t^*(1)\) 到 \(K\) 上的限制为有限且连续的紧子集组成）是 \(\mu\) -稠密的，且 §2, No.~3 的命题 4 蕴涵存在正测度的一个可和族 \((\mu_\alpha)_{\alpha \in A}\) ，其支集属于 \(\mathfrak{K}\) ，使得 \(\mu = \sum_{\alpha \in A} \mu_\alpha\) 。映射 \(\Lambda\) 是 \(\mu_\alpha\) -适宜的（对每个 \(\alpha \in A\) ）；令 \(\nu_\alpha = \int \lambda_t d\mu_\alpha(t)\) 。命题 1 表明 \(\nu = \sum_{\alpha \in A} \nu_\alpha\) ，而把关系式 (4) 应用于测度 \(\mu_\alpha\) 和函数 1，就表明 \(\nu_\alpha\) 是有界测度（因为 \(\lambda_t^*(1)\) 在 \(\operatorname{Supp}(\mu_\alpha)\) 上有界）。于是对测度 \(\mu_\alpha\) 写出公式 (6) ，并将第一个成员中的符号 \(\int^*\) 换成 \(\int^\bullet\) ，由 §1 的命题 7，这是合法的；则

\[
\int^ {\bullet} f (x) d \nu_ {\alpha} (x) \geqslant \int^ {\bullet} d \mu_ {\alpha} (t) \int^ {*} f (x) d \lambda_ {t} (x) = \int^ {\bullet} d \mu_ {\alpha} (t) \int^ {\bullet} f (x) d \lambda_ {t} (x)
\]

（最后一个等式是由于 \(\lambda_t\) 局部几乎处处有界，以及 § 1 的命题 7）。然后对 \(\alpha\) 求和即得 (8) 中的不等式（§2, No.~2, 命题 1）。

如果不作任何与 c) 类似的假设，不等式 (8) 可能不成立（习题 2）。

推论 1. --- 设 \(f\) 是一个 \(\geqslant 0\) 的函数，定义在 \(X\) 上，并设 \(H\) 是由 \(t \in T\) （使 \(f\) 非 \(\lambda_t\) -可忽略者）组成的集。

\begin{enumerate}
\def\labelenumi{\alph{enumi})}
\item
  若 \(f \nu\) -可忽略，则 \(H\) 是局部 \(\mu\) -可忽略的。
\item
  若 \(f\) 是 \(\nu\) -可忽略的且 \(\Lambda\) 是淡连续的，则 \(H\) 是 \(\mu\) -可忽略的。
\item
  若 \(f\) 是局部 \(\nu\) -可忽略的且 \(\lambda_t^\bullet(1) < +\infty\) 局部 \(\mu\) -几乎处处成立，则 \(H\) 是局部 \(\mu\) -可忽略的。
\end{enumerate}

推论 2. --- 设 f 是定义在 X 上的一个函数 \(\geqslant0\) ，为 \(\nu\) -可测且 \(\nu\) -适度的。于是由 \(t\in T\) （使 f 非 \(\lambda_{t}\) -适度者）组成的集是局部 \(\mu\) -可忽略的（甚至是 \(\mu\) -可忽略的，如果 \(\Lambda\) 是淡连续的）。

因为，\(f\) 是一列函数 \(f_{n} \geqslant 0\) 之和，其中 \(f_{n}\) 在一个紧集 \(K_{n}\) 之外为零（对 \(n \geqslant 1\) ），而 \(f_{0}\) 是 \(\nu\) -可忽略的（§1, No.~2, 命题 6）；于是 \(f_{0}\) 是 \(\lambda_{t}\) -可忽略的，但 \(t\) 构成局部 \(\mu\) -可忽略集（甚至是 \(\mu\) -可忽略集，如果 \(\Lambda\) 是淡连续的）者除外（推论 1），从而命题立即得出。

命题 4. --- 设 f 是定义在 X 上、取值于拓扑空间 G 中的一个 \(\nu\) -可测函数，并设 M 是由 \(t \in T\) （使 f 非 \(\lambda_{t}\) -可测者）组成的集。

\begin{enumerate}
\def\labelenumi{\alph{enumi})}
\item
  假设 \(f\) 在一个 \(\nu\) -适度子集（\(X\) 中的）的余集上是常数；则 \(M\) 是局部 \(\mu\) -可忽略的。
\item
  假设 \(f\) 在一个 \(\nu\) -适度子集（\(X\) 中的）的余集上是常数，且 \(\Lambda\) 是淡连续的；则 \(M\) 是 \(\mu\) -可忽略的。
\item
  假设 \(\lambda_t^\bullet(1) < +\infty\) 局部 \(\mu\) -几乎处处成立；则 M 是局部 \(\mu\) -可忽略的。
\end{enumerate}

先证 a)（相应地，b)）。由于每个 \(\nu\) -可积集都含于一个 \(\nu\) -可积开集，函数 f 在可数个 \(\nu\) -可积开集之并的余集 B 上是常数。存在 X - B 的一个分割，由一个 \(\nu\) -可忽略集 N 与一列紧集 \((\mathrm{K}_{n})\) 组成，使得 f 到每个 \(K_{n}\) 上的限制连续。设 S 是由 \(t \in T\) （使 N 非 \(\lambda_{t}\) -可忽略者）组成的集：由命题 3 的推论 1，S 是局部 \(\mu\) -可忽略的（相应地，\(\mu\) -可忽略的）。集合 \(K_{n}, B, N\) 对 X 上每个测度都可测，且 f 到其中每一个上的限制是 \(\lambda_{t}\) -可测的（对每个 \(t \notin S\) ）。因此函数 f 是 \(\lambda_{t}\) -可测的（对每个 \(t \notin S\) ）（Ch. IV, §5, No.~10, 命题 16）。

为建立 c)，我们沿用命题 3 证明中的记号；由于 f 是 \(\nu\) -可测的，它对每个测度 \(\nu_{\alpha} \leqslant \nu\) 都可测。而这些测度是有界的，从而是适度的，于是由 a) 推出，M 是局部 \(\mu_{\alpha}\) -可忽略的（对每个 \(\alpha \in A\) ）。这蕴涵 M 是局部 \(\mu\) -可忽略的（§2, No.~2, 命题 1 的推论 2）。

命题 5. --- 设 f 是定义在 X 上的一个 \(\nu\) -可测正数值函数，并设 N 是由 \(t \in T\) （使 f 不同时为 \(\lambda_{t}\) -可测且 \(\lambda_{t}\) -适度者）组成的集。

\begin{enumerate}
\def\labelenumi{\alph{enumi})}
\tightlist
\item
  假设 \(f\) 是 \(\nu\) -适度的。于是集 \(\mathbf{N}\) 是局部 \(\mu\) -可忽略的，函数 \(t \mapsto \int^{\bullet} f(x) d\lambda_t(x)\) 是 \(\mu\) -可测的，且
\end{enumerate}

\[
\int^ {\bullet} f (x) d \nu (x) = \int^ {\bullet} d \mu (t) \int^ {\bullet} f (x) d \lambda_ {t} (x).\tag{9}
\]

\begin{enumerate}
\def\labelenumi{\alph{enumi})}
\setcounter{enumi}{1}
\tightlist
\item
  假设 \(f\) 是 \(\nu\) -适度的，且 \(\Lambda\) 是淡连续的。于是集 \(N\) 是 \(\mu\) -可忽略的，函数 \(t \mapsto \int^{*} f(x) d\lambda_{t}(x)\) 是 \(\mu\) -可测且 \(\mu\) -适度的，且
\end{enumerate}

\[
\int^ {*} f (x) d \nu (x) = \int^ {*} d \mu (t) \int^ {*} f (x) d \lambda_ {t} (x).\tag{10}
\]

\begin{enumerate}
\def\labelenumi{\alph{enumi})}
\setcounter{enumi}{2}
\tightlist
\item
  假设 \(\lambda_t^\bullet(1) < +\infty\) 局部 \(\mu\) -几乎处处成立。于是集 N 是局部 \(\mu\) -可忽略的，函数 \(t \mapsto \int^\bullet f(x) d\lambda_t(x)\) 是 \(\mu\) -可测的，且 (9) 成立。
\end{enumerate}

先在 \(f\) 是 \(\nu\) -适度的假设下证明 a)（相应地，b)）。关于集 \(N\) 的断言已经建立（命题 4，以及命题 3 的推论 2）。由 §1, No.~2 的命题 6，我们只需在下述每个特殊情形中证明 a)（相应地，b)）：

\begin{enumerate}
\def\labelenumi{\arabic{enumi})}
\item
  函数 \(f\) 是 \(\nu\) -可忽略的。
\item
  存在一个紧集 \(\mathbf{K}\) ，使得 \(f\) 在 \(\mathbf{K}\) 之外为零且 \(f\) 到 \(\mathbf{K}\) 上的限制连续。
\end{enumerate}

特殊情形 1) 已经处理过（命题 3 的推论 1）。为处理第二个情形，用 G 表示包含 K 的一个 \(\nu\) -可积开集，用 M 表示 f 的一个常数上界，用 h 表示下半连续函数 \(M\varphi_{G}\) ，并用 g 表示函数 h-f.~由于函数 f 在 X 上是上半连续的，g 是下半连续且正的。此外，f,g,h 都是 \(\nu\) -可积的。

于是把公式 (4)（相应地，(5)）应用于下半连续函数 h 与 g。由减法可见，函数

\[
t \mapsto \int^ {\bullet} f (x) d \lambda_ {t} (x) \quad (\text { resp. } \int^ {*} f (x) d \lambda_ {t} (x))
\]

是 \(\mu\) -可测的，且公式 (9)（相应地，(10)）成立。最后，在 b) 的假设下，函数 \(t \mapsto \int^{*} f(x) d\lambda_{t}(x)\) 具有有限的上积分，因而确实是 \(\mu\) -适度的。

为证 c)，我们沿用命题 3 证明中的测度 \(\mu_{\alpha}\) 与 \(\nu_{\alpha}\) ；由于 f 是 \(\nu_{\alpha}\) -可测且 \(\nu_{\alpha}\) -适度的，断言 a) 蕴涵 \(t \mapsto \int^{\bullet} f(x) d\lambda_{t}(x)\) 是 \(\mu_{\alpha}\) -可测的，且

\[
\int^ {\bullet} f (x) d \nu_ {\alpha} (x) = \int^ {\bullet} d \mu_ {\alpha} (t) \int^ {\bullet} f (x) d \lambda_ {t} (x).
\]

只需再对 \(\alpha\) 求和，并应用 §2, No.~2 的命题 1 与命题 2。

如果不假设 \(f\) 是 \(\nu\) -适度的，并且不作 c) 中的假设，则关系式 (9) 可能不成立（习题 3）。

推论. --- 设 f 是定义在 X 上、取值于 Banach 空间 F 或 \(\overline{R}\) 中的函数，且它是 \(\nu\) -可测并 \(\nu\) -适度的。为使 f 是 \(\nu\) -可积的，必须且只须

\[
\int^ {\bullet} d \mu (t) \int^ {\bullet} | \mathbf {f} (x) | d \lambda_ {t} (x) <   + \infty .
\]

这由命题 5 及可积性判别法（Ch. IV, §5, No.~6, 定理 5）立即得出。

